\documentclass[12pt,a4paper]{article}


\usepackage{graphicx}



\textheight 240mm  %240/12
\textwidth  160mm   %160/12
\topmargin -10mm
\begin{document}

\title{Applications of Singular-Value Decomposition (SVD) }

\author{Alkiviadis G. Akritas\footnote{University of Thessaly, Department of
  Computer and Communication Engineering, GR-38221 Volos,
  Greece}\hspace{.1in} and\hspace{.1in}Gennadi I. Malaschonok\footnote{Tambov University,
  Department of Mathematics,Tambov, Russia}}

\date{Summer 2002}
\maketitle

\begin{abstract}
Let $A$ be an $m \times n$ matrix with $m \geq n$. Then one form
of the singular-value decomposition of $A$ is

$$ A = U^T \Sigma V,$$

where $U$ and $V$ are orthogonal and $\Sigma$ is square diagonal.
That is, $UU^T=I_{rank(A)},$ $VV^T=I_{rank(A)},$ $U$ is $rank(A)
\times m,$ $V$ is $rank(A) \times n$ and

$$ \Sigma = \left(\begin{array}{ccccc}  \sigma_1&0&\cdots&0&0\\
                                                                 0&\sigma_2&\cdots&0&0\\
                                                                 \vdots&\vdots&\ddots&\vdots&\vdots\\
                                                                 0&0&\cdots&\sigma_{rank(A)-1}&0\\
                                                                 0&0&\cdots&0&\sigma_{rank(A)}
\end{array}\right)  $$

is a $rank(A) \times rank(A)$ diagonal matrix.  In addition
$\sigma_1\geq \sigma_2\geq \cdots \geq \sigma_{rank(A)} >  0.$ The
$\sigma_i$'s are called the {\em singular values} of $A$ and their
number is equal to the rank of $A$.  The ratio $\sigma_1 \over
\sigma_{rank(A)} $ can be regarded as a condition number of the
matrix $A$.

It is easily verified that the singular-value decomposition can be
also written as

$$ A = U^T \Sigma V = \sum_{i=1}^{rank(A)} \sigma_i u_i^T v_i.$$

The matrix $ u_i^T v_i$ is the {\em outer product} of the i-th row
of $U$ with the corresponding row of $V$.  Note that each of these
matrices can be stored using only $m+n$ locations rather than $mn$
locations.


The singular value decomposition is over a hundred years old. For
the case of square matrices, it was discovered independently by
Beltrami in 1873 and Jordan in 1874. The technique was extended to
rectangular matrices by Eckart and Young in the 1930's and its use
as a computational tool dates back to the 1960's.  Gene Golub and
van Loan demonstrated its usefulness and feasibility in a wide
variety of applications.

Using both forms presented above---and following Jerry Uhl's
beautiful approach---we show how SVD can be used as a tool for
teaching Linear Algebra geometrically, and then apply it in
solving least-squares problems and in data compression.

\end{abstract}



\end{document}
